\documentclass[11pt]{article}
\usepackage{mathblatt}
\usepackage{multicol}


\begin{document}
\pfheftstil
\pfheftkopf{Daten \textperiodcentered{} Lösungen}{}
\begin{pfloesung}{}
\lz{1.}{19,2 < 35,3 < 52,2 < 85,1}{}{}
\end{pfloesung}
\begin{pfloesung}{}
\lz{2.}{200; 48; 4}{}{}
\end{pfloesung}
\begin{pfloesung}{}
\lz{3.}{360°; 180°; 90°; 36°; 3,6°}{}{}
\end{pfloesung}
\begin{pfloesung}{}
\lz{4.}{20 je Kästchen; 4,5 Kästchen}{}{}
\end{pfloesung}
\begin{pfloesung}{}
\lz{5.}{50 \%; 25 \%; 75 \%; $\tfrac{1}{15}$ = $\tfrac{20}{3}$ \% $\approx$ 6,7 \%}{}{}
\end{pfloesung}
\begin{pfloesung}{}
\lz{6.}{Modalwert}{}{}
\end{pfloesung}
\begin{pfloesung}{}
\lz{7.}{$22$}{Minimum $= 9$, Maximum }{}
\end{pfloesung}
\begin{pfloesung}{}
\lz{8.}{74 mm $-$ 52 mm = 22 mm}{}{}
\end{pfloesung}
\begin{pfloesung}{}
\lz{9.}{2 h 43 min (Mittwoch)}{}{}
\end{pfloesung}
\begin{pfloesung}{}
\lz{10.}{40 oder 46}{}{}
\end{pfloesung}
\begin{pfloesung}{}
\lz{11.}{wärmster Monat: August (15 °C); kältester Monat: Januar ($-$4 °C)}{}{}
\end{pfloesung}
\begin{pfloesung}{}
\lz{12.}{Minimum 154,2 ct; Maximum 169,2 ct}{}{}
\end{pfloesung}
\begin{pfloesung}{}
\lz{13.}{14,4 ct}{154,0 ct $-$ 139,6 ct}{}
\end{pfloesung}
\begin{pfloesung}{}
\lz{14.}{Minimum 400 (2021), Maximum 485 (2019); Mittel 451}{Summe 2255 : 5}{}
\end{pfloesung}
\begin{pfloesung}{}
\lz{15.}{0,4 m}{}{}
\end{pfloesung}
\begin{pfloesung}{}
\lz{16.}{65,9 mm}{85,1 $-$ 19,2}{}
\end{pfloesung}
\begin{pfloesung}{}
\lz{17.}{Minimum 290 Mio. (Arabisch); Maximum 1300 Mio. (Chinesisch); Spannweite 1010 Mio.}{1300 $-$ 290}{}
\end{pfloesung}
\begin{pfloesung}{}
\lz{18.}{280 $-$ 130 = 150; 210 + 210 : 3 = 280}{ein Drittel von 210 = 70}{}
\end{pfloesung}
\begin{pfloesung}{}
\lz{19.}{2,7 Mio.; 8,7 Mio.; $\tfrac{11}{79}$ $\approx$ 13,9 \%}{9,7 $-$ 7,0 = 2,7; Summe 78,3 Mio.; 78,3 : 9 = 8,7; 9,0 $-$ 7,9 = 1,1 Mio.}{}
\end{pfloesung}
\begin{pfloesung}{}
\lz{20.}{4 · 8,5 = 34 Treffer}{}{}
\end{pfloesung}
\begin{pfloesung}{}
\lz{21.}{36}{}{}
\end{pfloesung}
\begin{pfloesung}{}
\lz{22.}{$1{,}3$}{Tabelle: $0$: $2$; $1$: $4$; $2$: $3$; $3$: $1$; $\bar{x} = (0 + 4 + 6 + 3) : 10$}{}
\end{pfloesung}
\begin{pfloesung}{}
\lz{23.}{(154 + 170 + 162 + 158 + 156) : 5 = 800 : 5 = 160 cm}{}{}
\end{pfloesung}
\begin{pfloesung}{}
\lz{24.}{200 Besucher}{Summe 1400}{}
\end{pfloesung}
\begin{pfloesung}{}
\lz{25.}{$17$}{Mittel $= 216 : 8 = 27$; Spannweite $= 35 - 18$}{}
\end{pfloesung}
\begin{pfloesung}{}
\lz{26.}{Mittelwert 0 °C; Unterschied 24 K}{}{}
\end{pfloesung}
\begin{pfloesung}{}
\lz{27.}{1,7 m}{}{}
\end{pfloesung}
\begin{pfloesung}{}
\lz{28.}{4,00 m}{Summe 16,00 m}{}
\end{pfloesung}
\begin{pfloesung}{}
\lz{29.}{327 mm : 6 = 54,5 mm}{}{}
\end{pfloesung}
\begin{pfloesung}{}
\lz{30.}{$2{,}92$}{$n = 4 + 6 + 7 + 5 + 2 + 1 = 25$; $\bar{x} = (1 \cdot 4 + 2 \cdot 6 + 3 \cdot 7 + 4 \cdot 5 + 5 \cdot 2 + 6 \cdot 1) : 25 = 73 : 25$}{}
\end{pfloesung}
\begin{pfloesung}{}
\lz{31.}{Mittelwert: 92,4 : 5 = 18,48 $\approx$ 18,5 °C – Aussage stimmt}{}{}
\end{pfloesung}
\begin{pfloesung}{}
\lz{32.}{680 $\tfrac{353}{17}$ $\approx$ 40 000 Zuschauer (40 021)}{680 353 : 17 $\approx$ 40 020,8}{}
\end{pfloesung}
\begin{pfloesung}{}
\lz{33.}{53 : 20 = 2,65; Klasse 9 (2,65 < 2,8)}{}{}
\end{pfloesung}
\begin{pfloesung}{}
\lz{34.}{581,7 : 12 = 48,475 $\approx$ 48,5 mm; April liegt $\tfrac{293}{485}$ $\approx$ 60,4 \% unter dem Durchschnitt}{Differenz 48,5 $-$ 19,2 = 29,3 mm; 29,3 : 48,5 = $\tfrac{293}{485}$ $\approx$ 0,604}{}
\end{pfloesung}
\begin{pfloesung}{}
\lz{35.}{Anteil in \%}{}{}
\end{pfloesung}
\begin{pfloesung}{}
\lz{36.}{Anteil in \%}{}{}
\end{pfloesung}
\begin{pfloesung}{}
\lz{37.}{$3$ cm}{$30$ mm}{}
\end{pfloesung}
\begin{pfloesung}{}
\lz{38.}{$36^\circ$}{ein Zehntel von $360^\circ$}{}
\end{pfloesung}
\begin{pfloesung}{}
\lz{39.}{Fußball 90°, Basketball 54°}{}{}
\end{pfloesung}
\begin{pfloesung}{}
\lz{40.}{Sektor $72^\circ$ ab der Nullmarke, beschriftet}{}{}
\end{pfloesung}
\begin{pfloesung}{}
\lz{41.}{216°, 108°, 36°}{}{}
\end{pfloesung}
\begin{pfloesung}{}
\lz{42.}{10 \%; ab 12 Uhr im Uhrzeigersinn: Getreide, Flüssigkeiten, Schrott, Sonstiges, Container}{100 $-$ 17 $-$ 66 $-$ 6 $-$ 1 = 10}{}
\end{pfloesung}
\begin{pfloesung}{}
\lz{43.}{$36$}{Winkel $= 180$; $144$; in Grad}{}
\end{pfloesung}
\begin{pfloesung}{}
\lz{44.}{der Rasen}{Haus $= 20\,\%$ , Rasen $= 50\,\%$ , Beete $= 15\,\%$ , Wege $= 15\,\%$; $72^\circ$; $180^\circ$; $54^\circ$; $54^\circ$}{}
\end{pfloesung}
\begin{pfloesung}{}
\lz{45.}{180°, 126°, 27°, 27°}{}{}
\end{pfloesung}
\begin{pfloesung}{}
\lz{46.}{grün 108°, gelb 72°, blau 45°, rot 135°}{}{}
\end{pfloesung}
\begin{pfloesung}{}
\lz{47.}{$25\,\%$}{$100 - 45 - 30$}{}
\end{pfloesung}
\begin{pfloesung}{}
\lz{48.}{größter Sektor Rad, dann Bus, dann Auto, kleinster zu Fuß}{}{}
\end{pfloesung}
\begin{pfloesung}{}
\lz{49.}{A 22,5°; B 180°; C 45°; D 90°; E 22,5°}{}{}
\end{pfloesung}
\begin{pfloesung}{}
\lz{50.}{rechts „einmal pro Woche“, unten „einmal im Monat“, links „mehrmals pro Woche“; 57,6°}{0,16 · 360°}{}
\end{pfloesung}
\begin{pfloesung}{}
\lz{51.}{Sektor links unten vor Gartenpflege ($\approx$ 43°) beschriften; $\tfrac{1080}{127}$° $\approx$ 8,5°}{60 : 508 · 360° = $\tfrac{5400}{127}$° $\approx$ 42,5°; 12 : 508 = $\tfrac{3}{127}$ $\approx$ 0,0236}{}
\end{pfloesung}
\begin{pfloesung}{}
\lz{52.}{viertes oder weiteres Kind; 40,32° $\approx$ 40°; Sektor 40° „drittes Kind“}{5 \% $\mathrel{\widehat{=}}$ 18°; 0,112 · 360° = 40,32° $\approx$ 40,3°}{}
\end{pfloesung}
\begin{pfloesung}{}
\lz{53.}{ja, bei null}{}{}
\end{pfloesung}
\begin{pfloesung}{}
\lz{54.}{$0{,}1$ je Linie}{$0{,}5 : 5$}{}
\end{pfloesung}
\begin{pfloesung}{}
\lz{55.}{Stäbe bei 300, 120, 60 und 20}{}{}
\end{pfloesung}
\begin{pfloesung}{}
\lz{56.}{Säule für Tag 7 bis 20 km}{}{}
\end{pfloesung}
\begin{pfloesung}{}
\lz{57.}{$40$ je Kästchen}{$240 : 6$}{}
\end{pfloesung}
\begin{pfloesung}{}
\lz{58.}{Säulen $18$, $24$, $6$, $12$}{Einteilung in Sechserschritten bis $30$}{}
\end{pfloesung}
\begin{pfloesung}{}
\lz{59.}{Punkte bei $1$, $3$, $7$, $12$, $16$, verbunden zu einer Linie}{}{}
\end{pfloesung}
\begin{pfloesung}{}
\lz{60.}{Balken bei Do bis 180}{}{}
\end{pfloesung}
\begin{pfloesung}{}
\lz{61.}{$120$}{Handball $= 280$; Tennis ; die Säule mit dem Fragezeichen; zu zeichnen}{}
\end{pfloesung}
\begin{pfloesung}{}
\lz{62.}{Kopf $\tfrac{12}{25}$ = 48 \%, Zahl $\tfrac{13}{25}$ = 52 \%}{}{}
\end{pfloesung}
\begin{pfloesung}{}
\lz{63.}{20 je Gitterlinie (0 bis 200); E und D Kartoffel- und Gemüsegerichte, C Pizza, B Schnitzel, A Fischgerichte}{längste Balken 160 $\Rightarrow$ 20 je Gitterlinie}{}
\end{pfloesung}
\begin{pfloesung}{}
\lz{64.}{$40$ Besucher}{}{}
\end{pfloesung}
\begin{pfloesung}{}
\lz{65.}{absolut: 4 $\to$ 75, 5 $\to$ 114; relativ: 2 $\to$ 0,116, 5 $\to$ 0,228}{}{}
\end{pfloesung}
\begin{pfloesung}{}
\lz{66.}{1 Kästchen = 50 Mio. (z. B. 100, 200, … je zwei Kästchen); Hindi; Säule Arabisch 290 : 50 = 5,8 Kästchen hoch}{389 Mio. $\mathrel{\widehat{=}}$ 7,8 Kästchen $\Rightarrow$ 50 Mio. je Kästchen; 10,5 · 50 = 525}{}
\end{pfloesung}
\begin{pfloesung}{}
\lz{67.}{bei null}{}{}
\end{pfloesung}
\begin{pfloesung}{}
\lz{68.}{nicht bei null, Startwert $80$}{}{}
\end{pfloesung}
\begin{pfloesung}{}
\lz{69.}{die Aussage stimmt}{Jeder zwanzigste ist $\tfrac{1}{20} = 5\,\%$}{}
\end{pfloesung}
\begin{pfloesung}{}
\lz{70.}{mehr als die Hälfte: 13 > 10,5}{}{}
\end{pfloesung}
\begin{pfloesung}{}
\lz{71.}{Nein}{15 Cent}{}
\end{pfloesung}
\begin{pfloesung}{}
\lz{72.}{die Aussage stimmt}{Ein Drittel von $240$ ist $= 80$, $240 + 80 = 320$}{}
\end{pfloesung}
\begin{pfloesung}{}
\lz{73.}{aber die Fortschreibung ist unzulässig: nichts im Diagramm sagt, dass der Rückgang so weitergeht}{Rechnerisch $320 : 20 = 16$ Jahre}{}
\end{pfloesung}
\begin{pfloesung}{}
\lz{74.}{$52$, weit unter der Hälfte}{Teil 1 falsch: 2021 ist die Zahl gestiegen; Teil 2 falsch: von $348$ auf $296$ ist ein Rückgang um ; $348$ auf $361$}{}
\end{pfloesung}
\begin{pfloesung}{}
\lz{75.}{z. B. „extrem abnehmen“ falsch: von 44 \% auf 36 \%, abgeschnittene Achse übertreibt}{8 Prozentpunkte in 4 Jahren}{}
\end{pfloesung}
\begin{pfloesung}{}
\lz{76.}{$20$ je Jahr, bis $0$ wären es über $17$ Jahre, und ein Trend gilt nicht beliebig lange}{Teil 1 falsch: 2022 stieg die Zahl; Teil 2 unbelegt: im Schnitt sinkt die Zahl um ; $380$ auf $395$}{}
\end{pfloesung}
\begin{pfloesung}{}
\lz{77.}{Ja. 17 von 20 = 85 \% > 75 \%.}{}{}
\end{pfloesung}
\begin{pfloesung}{}
\lz{78.}{$\tfrac{1}{36} \approx 2{,}8\,\%$}{Sichtbar sind nur $2$ und $4$ – das verdoppelt sich; der Wert steigt von $72$ auf $74$, also um $= 2 : 72$}{}
\end{pfloesung}
\begin{pfloesung}{}
\lz{79.}{falsch}{2018 $\to$ 2019 Anstieg 432 auf 485 | 2018 $\to$ 2021 Rückgang nur um 32 ($\tfrac{2}{27}$ $\approx$ 7,4 \%), nicht um die Hälfte; 432 $-$ 400 = 32; Hälfte von 432 = 216}{}
\end{pfloesung}
\begin{pfloesung}{}
\lz{80.}{Aussage 1 falsch (55 \% > 50 \%); Aussage 2 falsch (jede 15. = $\tfrac{1}{15}$ $\approx$ 6,7 \%, tatsächlich 15 \%)}{}{}
\end{pfloesung}
\begin{pfloesung}{}
\lz{81.}{5 °C}{}{}
\end{pfloesung}
\begin{pfloesung}{}
\lz{82.}{19 °C}{sortiert 11, 17, 18, 20, 21, 21; Mitte zwischen 18 und 20}{}
\end{pfloesung}
\begin{pfloesung}{}
\lz{83.}{1~Mittelwert, 2~Minimum, 3~Median, 4~Mittelwert, 5~Minimum, 6~Median}{}{}
\end{pfloesung}
\begin{pfloesung}{}
\lz{84.}{18 °C}{Summe der sechs Werte 122; 7 · 20 = 140}{}
\end{pfloesung}
\begin{pfloesung}{}
\lz{85.}{Mehr als die Hälfte der Spiele hat 0 Punkte; der Median zeigt nicht, dass es auch Spiele mit 1 oder 2 Punkten gab.}{}{}
\end{pfloesung}
\begin{pfloesung}{}
\lz{86.}{$4$}{Aussage 2 – der häufigste Wert ist $= 3$ , nicht $5$; Aussage 1 stimmt: $6 - 2$; viermal; dreimal}{}
\end{pfloesung}
\begin{pfloesung}{}
\lz{87.}{Er steigt, weil der neue Wert über dem bisherigen Mittelwert liegt}{}{}
\end{pfloesung}
\begin{pfloesung}{}
\lz{88.}{$11$}{Mittelwert $= 95 : 5 = 19$; Median }{}
\end{pfloesung}
\begin{pfloesung}{}
\lz{89.}{Aussage 2 falsch}{Modalwert ist 2 Jahre (fünfmal), Spannweite 5 $-$ 1 = 4 stimmt; Häufigkeiten: 1 ×1, 2 ×5, 4 ×3, 5 ×1}{}
\end{pfloesung}
\begin{pfloesung}{}
\lz{90.}{48 Jahre; bleibt 48 Jahre, denn 66 + 30 = 96 = 2 · 48 (Abweichungen +18 und $-$18 heben sich auf)}{Summe 528 : 11; neue Summe 432 : 9}{}
\end{pfloesung}
\begin{pfloesung}{}
\lz{91.}{$20\,\%$}{$72 : 360 = 0{,}2$}{}
\end{pfloesung}
\begin{pfloesung}{}
\lz{92.}{Streifen: Lebensmittel 1/4, Miete 3/10, Sonstiges $\tfrac{45}{100}$ der Länge}{}{}
\end{pfloesung}
\begin{pfloesung}{}
\lz{93.}{Streifen mit 1 cm, 1,5 cm, 5,5 cm, 2 cm, beschriftet}{10 cm $\mathrel{\widehat{=}}$ 100 \% $\Rightarrow$ 1 mm je 1 \%}{}
\end{pfloesung}
\begin{pfloesung}{}
\lz{94.}{$10$ Prozentpunkte: Erdbeere stieg um $10$, Vanille sank um $10$}{}{}
\end{pfloesung}
\begin{pfloesung}{}
\lz{95.}{B: 11 \%; Streifen 0: 4,1 cm, A: 4,3 cm, B: 1,1 cm, AB: 0,5 cm}{}{}
\end{pfloesung}
\begin{pfloesung}{}
\lz{96.}{$\tfrac{5}{11}$ $\approx$ 45,5 \%; Sektor $\tfrac{1800}{11}$° $\approx$ 164° (Rest $\tfrac{2160}{11}$° $\approx$ 196°)}{5 von 11 Personen; $\tfrac{5}{11}$ · 360° = $\tfrac{1800}{11}$° $\approx$ 163,6°}{}
\end{pfloesung}
\begin{pfloesung}{}
\lz{97.}{Nur dann ist die Säulenhöhe zum Wert proportional}{}{}
\end{pfloesung}
\begin{pfloesung}{}
\lz{98.}{$2$ Grad, gut $10\,\%$}{Die Achse beginnt bei $18$ statt bei $0$; sichtbar sind nur $1$ und $3$ Grad über dem Achsenanfang; Diese Teile verhalten sich wie $1 : 3$, die Temperatur stieg aber nur um }{}
\end{pfloesung}
\begin{pfloesung}{}
\lz{99.}{$3$, die Werte selbst unterscheiden sich nur um knapp $5\,\%$}{Ab $82$ zeigen die Säulen nur $2$, $4$ und $6$; das Verhältnis der sichtbaren Teile ist }{}
\end{pfloesung}
\begin{pfloesung}{}
\lz{100.}{Säulen bei $1\,020$, $1\,040$, $1\,060$ – fast gleich hoch}{knapp $4\,\%$}{}
\end{pfloesung}
\begin{pfloesung}{}
\lz{101.}{y-Achse beginnt bei 1060 € statt bei 0 $\Rightarrow$ kleine Zuwächse wirken groß}{Zuwachs je Jahr nur wenige Euro}{}
\end{pfloesung}
\begin{pfloesung}{\textbf{Prüfstein}}
\lz{a)}{0,08 € pro Liter}{1,85 $-$ 1,77}{}
\lz{b)}{wahr}{Mittel Mo–Fr = 8,95 : 5 = 1,79 € < 1,80 €; Summe 8,95}{}
\lz{c)}{$\tfrac{5}{179}$ $\approx$ 2,8 \%}{Differenz 0,05 €; 0,05 : 1,79 = $\tfrac{5}{179}$ $\approx$ 0,0279}{}
\lz{d)}{Säule So bis 1,85 (höchste Säule)}{Skala 1,75 bis 1,90, Schritt 0,01}{}
\lz{e)}{Achse beginnt bei 1,75 € statt bei 0 $\Rightarrow$ Säulen zeigen nur die Differenz zu 1,75 € (Do 0,04, Fr 0,09) $\Rightarrow$ Verdopplung der Säulenhöhe, Preis nur +2,8 \%}{Do 1,79 $-$ 1,75 = 0,04; Fr 1,84 $-$ 1,75 = 0,09}{}
\end{pfloesung}
\end{document}